\documentclass[11pt]{article}
\usepackage[margin=1in]{geometry}
\usepackage{graphicx}
\usepackage{booktabs}
\usepackage{amsmath}
\usepackage{amssymb}
\usepackage{xcolor}
\usepackage[hidelinks]{hyperref}
\usepackage{caption}
\captionsetup{font=small}

\title{\vspace{-2.5em}Sparks in detector~7, and how they compare to detector~3\\
\large (June cosmic bench --- \texttt{g\_det7\_long})}
\author{MX17 cosmic-bench QA}
\date{\today}

\begin{document}
\maketitle
\vspace{-1.5em}

\section{Purpose}
The det3 spark study (companion note \texttt{det3\_spark\_analysis}) found sparks that are
Poisson in time, predominantly muon-induced, edge-localised, and neither a single flash
nor a propagating discharge. det7 is the most spark-prone detector of the June fleet
(38.9\,\% of firing events spark, vs 9.1\,\% for det3) and runs at a different operating
point (drift 700~V, resist past optimum). We repeat the identical analysis on
\texttt{mx17\_det6\_det7\_overnight\_6-26-26 / long\_run} (det7 $=$ FEU6/8, 3.84~h,
\textbf{22\,247 sparks}) to ask: is it the \emph{same} phenomenon, just more of it?

\section{Answer: yes --- the same story, amplified}
Every det3 conclusion reproduces in det7, scaled up by the higher operating gain
(Table~\ref{tab:cmp}). Sparks are muon-seeded avalanche breakdowns whose probability and
spatial/temporal extent grow as the detector is pushed closer to breakdown, initiated at
the same edge / hot-strip weak spots, developing as a spatially-extended,
temporally-spread event with no propagating front.

\begin{table}[h]\centering\small
\begin{tabular}{lccl}
\toprule
 & \textbf{det3} (\texttt{g\_det3\_wknd}) & \textbf{det7} (\texttt{g\_det7\_long}) & reading\\
\midrule
spark rate (of firing events) & 9.1\,\% & 38.9\,\% & det7 deep in sparking regime\\
mean rate & 0.33~Hz & 1.61~Hz & higher Poisson mean\\
random in time? & Poisson & Poisson & neither bursts (no HV instab.)\\[2pt]
spark rate, muon \emph{crosses} det & 4.9\,\% & 31.0\,\% & muon-induced\\
spark rate, muon \emph{misses} ($>$30mm) & 1.2\,\% & 5.3\,\% & (stronger in det7)\\
\textbf{crossing/miss enhancement} & \textbf{4.1$\times$} & \textbf{5.9$\times$} & muon-dominated in both\\
spontaneous floor (miss rate) & 1.2\,\% & 5.3\,\% & floor rises with gain\\
sparks with no clean M3 track & 60\,\% & 47\,\% & DAQ cross-talk in both\\[2pt]
strips/spark (X, Y) & 40, 53 & 64, 76 & bigger discharges; Y$>$X in both\\
spatial span (median) & $\sim$200~mm (half) & $\sim$360~mm (whole) & det7 fills the plane\\
occupancy peaks at & edges $+$ hot strips & edges $+$ hot strips & same weak spots\\[2pt]
hit-time spread (median) & 253~ns & 334~ns & longer in det7\\
fraction $\sigma_t>300$~ns & 24\,\% & 66\,\% & det7 rings longer\\
position--time corr $|r|$ & 0.12 & 0.15 & no ``walk'' in either\\
\bottomrule
\end{tabular}
\caption{det3 vs det7 --- identical phenomenology, det7 simply further into breakdown.}
\label{tab:cmp}
\end{table}

\section{Q1 --- Time: Poisson, faster}
det7 sparks at a flat $1.61$~Hz with no bursts; inter-spark intervals are exponential and
counts-per-bin have $\mathrm{std}=\sqrt{\mathrm{mean}}$ (Fig.~\ref{fig:time}) --- Poisson,
exactly as det3, just a $5\times$ higher mean. The elevated rate is therefore \emph{not} a
new instability; it is the same random process with more muons reaching breakdown. (Mild
consecutive-event clustering, $P(\text{spark}\,|\,\text{prev spark})=39\,\%$ vs $32\,\%$
baseline, matches det3's small excess.)

\begin{figure}[h]\centering
\includegraphics[width=\textwidth]{fig_time.png}
\caption{det7 time structure: flat rate (a), exponential intervals (b), Poisson
counts/bin (c).}
\label{fig:time}
\end{figure}

\section{Q2 --- Muon-induced, even more so}
Conditioning on a clean single M3 muon, a track crossing det7 sparks at \textbf{31.0\,\%}
versus \textbf{5.3\,\%} for one that clearly misses --- a \textbf{5.9$\times$} enhancement
(Fig.~\ref{fig:muon}a), stronger than det3's 4.1$\times$. So $\sim$83\,\% of in-detector
sparks are muon-induced. The spontaneous floor is higher than det3 (5.3\,\% vs 1.2\,\%),
consistent with running closer to breakdown where even a stray electron can trigger a
discharge. The discharge centroid is again central and compact (Fig.~\ref{fig:muon}c). The
M3 crossing map (Fig.~\ref{fig:muon}b) shows more structure than det3's --- but this tracks
the overall M3 track density / telescope acceptance, not a det7 hot-spot; the
detector-frame centroid is the clean spatial measure. The ``no clean track'' population
persists (47\,\% of sparks vs 35\,\% baseline; same 1.3$\times$ over-representation and
same discharge size), the DAQ common-mode cross-talk flagged for det3.

\paragraph{Shower hypothesis ruled out (as in det3).} If the no-track sparks were
multi-particle showers they would leave M3 with $\ge2$ track candidates; instead only
\textbf{0.5\,\%} of det7 no-track sparks do (they are $62\,\%$ M3-blank, $38\,\%$ single
degraded track), and their discharge morphology is \emph{identical} to muon-tracked sparks
(median 122 vs 125 strips, span 389 vs 392~mm, spread 335 vs 333~ns; Fig.~\ref{fig:shower}).
A spark nearly doubles M3's no-track rate ($18\to29\,\%$ blank) --- cross-talk, not showers.
This reproduces the det3 result on a $2.5\times$ larger spark sample.

\begin{figure}[h]\centering
\includegraphics[width=\textwidth]{fig_muon.png}
\caption{det7 muon dependence: (a) crossing muons spark $5.9\times$ more than missing
ones; (b) M3 crossing positions; (c) central discharge centroid.}
\label{fig:muon}
\end{figure}

\begin{figure}[h]\centering
\includegraphics[width=0.86\textwidth]{fig_shower_test.png}
\caption{det7 shower test: (a) M3 candidate-track count by spark class; (b) the
$\ge2$-track shower fraction stays at the $\sim$0.5\,\% level; (c) no-track and
muon-tracked sparks have identical discharge size --- showers ruled out.}
\label{fig:shower}
\end{figure}

\section{Q3 --- Same places: edges and hot strips}
det7 strip occupancy peaks at the plane \textbf{edges} (both planes; strongest at low-$X$
and high-$Y$) and on a few \textbf{hot strips} --- most dramatically an X/FEU6 channel near
$45$~mm that fires in a large fraction of all sparks (Fig.~\ref{fig:places}a), plus fixed
Y/FEU8 spikes (Fig.~\ref{fig:places}b). This is the same edge$+$hot-strip signature as det3,
more pronounced. Discharges are larger (median 64 X / 76 Y strips) and now span essentially
the \emph{whole} plane ($\sim$360~mm of 399~mm), versus half the plane in det3 --- the
higher gain lets a discharge propagate across the full active area. The Y/FEU8 plane again
discharges harder than X, consistent with det7's known FEU8 saturation band. Only 14\,\% of
spark hits saturate the ADC.

\begin{figure}[h]\centering
\includegraphics[width=\textwidth]{fig_places.png}
\caption{det7 locations: (a,b) edge peaks $+$ hot strips in X and Y; (c) discharge size ---
larger than det3, Y harder than X.}
\label{fig:places}
\end{figure}

The det7 spark hitmaps (Fig.~\ref{fig:hitmap}) show the same fixed edge$+$hot-strip
\textbf{grid} footprint as det3 (with an even more prominent X/FEU6 hot strip), over a
central discharge centre --- a pattern locked to detector geometry, not the muon flux, and
so a shower veto in itself. The detector-side track check holds too: the widest contiguous
fired-strip cluster is median \textbf{20~mm} (Fig.~\ref{fig:clusters}c), far wider than a
$\sim$mm particle track, so the fired strips are discharge blobs, not tight shower tracks.

\begin{figure}[h]\centering
\includegraphics[width=\textwidth]{fig_hitmap.png}
\caption{det7 spark hitmaps: (a) discharge centre; (b) footprint grid (edges $+$ hot
strips); (c) per-plane occupancy.}
\label{fig:hitmap}
\end{figure}

\begin{figure}[h]\centering
\includegraphics[width=\textwidth]{fig_clusters.png}
\caption{det7 detector-side shower check: (a) clusters/spark; (b) fill fraction; (c)
widest contiguous cluster (median 20~mm $\gg$ a particle track).}
\label{fig:clusters}
\end{figure}

\section{Q4 --- Still no ``walk'', but longer-lived}
det7 sparks are, like det3, neither simultaneous ($0.8\,\%$ within 100~ns) nor propagating:
the position--time correlation is $|r|\approx0.15$ and the spark-centred
$\Delta\text{pos}$--$\Delta\text{time}$ map (Fig.~\ref{fig:walk}b) is again a cross with no
tilt. The difference is duration: the median per-spark time spread is $334$~ns (vs det3's
253~ns) and \textbf{66\,\%} of det7 sparks spread beyond $300$~ns (vs 24\,\% for det3) ---
the discharge rings/afterpulses for longer, filling much of the readout window
(Fig.~\ref{fig:walk}a,c). Same mechanism (extended discharge $+$ afterpulsing, timing
independent of position), longer tail.

\begin{figure}[h]\centering
\includegraphics[width=\textwidth]{fig_walk.png}
\caption{det7 time structure: (a) longer per-spark spread than det3; (b) no tilt (no
propagating front); (c) example sparks across the whole plane, spread in time.}
\label{fig:walk}
\end{figure}

\section{Takeaways}
\begin{itemize}\setlength\itemsep{0.15em}
  \item det7 and det3 spark by the \textbf{same mechanism}; det7 is simply operated further
        into breakdown, which raises the rate (0.33$\to$1.61~Hz), the muon enhancement
        (4.1$\to$5.9$\times$), the spontaneous floor (1.2$\to$5.3\,\%), the discharge size
        (half$\to$whole plane) and its duration (253$\to$334~ns).
  \item The lever is the \textbf{operating gain}: det7 needs a lower resist HV
        (the HV-scan optimum was $\sim$440~V; this run sat above it). Lowering it cuts the
        rate but not the muon-seeded nature.
  \item The \textbf{edge / hot-strip} weak spots and the \textbf{DAQ common-mode cross-talk}
        (spark $\to$ M3 track lost) reproduce on a second detector, so both are systemic,
        not det3 quirks --- worth fixing at the detector-edge/guard-ring and DAQ level.
\end{itemize}

\vspace{0.4em}
\footnotesize\noindent\textit{Reproduce:} \texttt{extract\_sparks.py g\_det7\_long} then
\texttt{make\_plots.py g\_det7\_long}; numbers in \texttt{spark\_meta.json}. Identical
pipeline to \texttt{det3\_spark\_analysis}.

\end{document}
